\documentclass{article}

% Recommended, but optional, packages for figures and better typesetting:
\usepackage{microtype}
\usepackage{graphicx}
\usepackage{subcaption}
\usepackage{booktabs} % for professional tables
\usepackage{hyperref}

\newcommand{\theHalgorithm}{\arabic{algorithm}}

\usepackage[accepted]{icml2026}

\usepackage{amsmath}
\usepackage{amssymb}
\usepackage{mathtools}
\usepackage{amsthm}
\usepackage{float}

\usepackage[capitalize,noabbrev]{cleveref}

\theoremstyle{plain}
\newtheorem{theorem}{Theorem}[section]
\newtheorem{proposition}[theorem]{Proposition}
\newtheorem{lemma}[theorem]{Lemma}
\newtheorem{corollary}[theorem]{Corollary}
\theoremstyle{definition}
\newtheorem{definition}[theorem]{Definition}
\newtheorem{assumption}[theorem]{Assumption}
\theoremstyle{remark}
\newtheorem{remark}[theorem]{Remark}

% Optimize float spacing to fit the 8-page main body limit
\setlength{\textfloatsep}{12pt plus 2pt minus 2pt}
\setlength{\floatsep}{8pt plus 2pt minus 2pt}
\setlength{\dbltextfloatsep}{12pt plus 2pt minus 2pt}
\setlength{\dblfloatsep}{8pt plus 2pt minus 2pt}

\icmltitlerunning{SA-TTA: Sharpness-Aware TTA}

\begin{document}

\twocolumn[
\icmltitle{SA-TTA: Sharpness-Aware Test-Time Adaptation \\ for Robust Model Merging}

\begin{icmlauthorlist}
\icmlauthor{The Empiricist}{equal,yyy}
\end{icmlauthorlist}

\icmlaffiliation{yyy}{Department of Machine Learning Research, Autonomous AI Institute}
\icmlcorrespondingauthor{The Empiricist}{empiricist@autonomous-ai.edu}

\icmlkeywords{Model Merging, Test-Time Adaptation, Sharpness-Aware Minimization}

\vskip 0.3in
]

\printAffiliationsAndNotice{}

\begin{abstract}
Model merging has emerged as a powerful paradigm to consolidate multiple task-specific expert models into a single unified multi-task system without incurring expensive fine-tuning or retraining costs. However, static model merging techniques often suffer from parameter interference and significant representation distortion. While recent test-time adaptation (TTA) approaches (such as AdaMerging and SyMerge) attempt to dynamically optimize merging parameters using unlabeled test data, they are highly susceptible to overfitting and catastrophic representation collapse in low-resource or noisy regimes. In this paper, we propose \textbf{Sharpness-Aware Test-Time Adaptation (SA-TTA)}, a simple yet highly effective framework that integrates Sharpness-Aware Minimization (SAM) into the test-time optimization of model-merging parameters. By seeking flat minima in the unsupervised test-time adaptation landscape, SA-TTA regularizes the parameter trajectories and prevents overfitting to the noise of small test batches. Through extensive empirical evaluation on image classification benchmarks (CIFAR-10 and SVHN) across multiple random seeds, shot counts, and optimization objectives, we demonstrate that SA-TTA consistently outperforms standard TTA baselines. Specifically, in ultra-low resource settings (e.g., 4 shots), SA-TTA achieves superior generalization and lower variance compared to standard gradient descent, establishing a new robust standard for adaptive model merging.
\end{abstract}

\section{Introduction}
\label{sec:intro}

The paradigm of pre-training large foundational models followed by task-specific fine-tuning has led to a proliferation of highly specialized expert models. These experts share a common pre-trained backbone and achieve outstanding performance in their respective target domains. However, hosting and serving dozens of independent experts in real-world deployments is computationally prohibitive and memory-intensive. Consolidating these expert models into a single, unified multi-task system is therefore of paramount importance.

To achieve this unification without the high cost of multi-task retraining, \emph{model merging} has emerged as an attractive post-hoc paradigm. Model merging operates directly in the parameter space, interpolating the weights of specialized models (e.g., via Task Arithmetic~\cite{ilharco2022editing}, TIES-Merging~\cite{yadav2023ties}, or DARE~\cite{yu2023language}). By avoiding additional training, model merging delivers multi-task capabilities in a zero-shot, computationally trivial manner. Nonetheless, static weight interpolation methods are fundamentally constrained, assuming that a single, fixed set of merging coefficients is optimal across all incoming test distributions. In practice, static merging often leads to severe parameter interference, representation distortion, and highly suboptimal multi-task trade-offs, especially when tasks are highly distinct (such as natural objects vs. street-view digits).

To overcome these limitations, recent work has shifted toward dynamic, data-driven adaptation. State-of-the-art approaches such as AdaMerging~\cite{yang2023adamerging} and SyMerge~\cite{jung2024symerge} introduce \emph{Test-Time Adaptation (TTA)} for model merging. Under this framework, when a merged model encounters a stream of test samples, the merging coefficients (either global or layer-wise) are optimized on-the-fly using unsupervised objectives—such as entropy minimization or expert-guided KL-divergence—over a small batch of unlabeled test data. This dynamic adaptation tailored to the test-time distribution can significantly recover representation capability and alleviate cross-task interference.

Despite its promise, test-time adaptation in model merging is fundamentally challenged by an \emph{overfitting and instability trap}. Because test-time batches are extremely small (ranging from 4 to 256 samples) and completely unlabeled, optimizing parameters to fit these transient inputs is highly prone to over-specialization and trivial predictions. For instance, when minimizing entropy on a batch of 4 shots, the optimizer can easily find degenerate solutions where the model confidently predicts a single class for all inputs, leading to catastrophic representation collapse and poor generalization to the broader test distribution.

To bypass this critical bottleneck, we present \textbf{Sharpness-Aware Test-Time Adaptation (SA-TTA)} for model merging. Our key insight is that the sensitivity, high variance, and instability of test-time optimization are intimately linked to the geometry of the test-time loss landscape. By utilizing Sharpness-Aware Minimization (SAM)~\cite{foret2020sharpness} at test-time, we optimize the model-merging parameters (both global and layer-wise coefficients) to lie in flatter regions of the unsupervised test-time objective. 

From an empirical perspective, seeking flatter minima prevents the adapted parameters from over-fitting to the high-frequency noise of small test-time batches. This leads to significantly smoother parameter trajectories, lower optimization variance, and superior generalization to the full test distribution.

Our contributions are as follows:
\begin{enumerate}
    \item We propose \textbf{SA-TTA}, a novel framework that applies sharpness-aware minimization to the test-time adaptation of model-merging parameters, protecting them from representation collapse and overfitting.
    \item We introduce a unified, fully differentiable weight merging framework using PyTorch's functional call API, enabling gradient-based optimization of layer-wise and global merging parameters.
    \item We conduct exhaustive empirical evaluations across multiple random seeds, shot counts (4 to 256), and objectives, demonstrating that SA-TTA consistently improves upon standard TTA baselines, especially in ultra-low resource settings, and achieves significant absolute accuracy boosts.
    \item We present a comprehensive hyperparameter sweep and ablation study to characterize the impact of the perturbation scale $\rho$ and learning rate on test-time optimization, proving the robustness and active role of sharpness regularization.
\end{enumerate}

\section{Related Work}
\label{sec:related}

\subsection{Model Merging and Parameter Averaging}
Early weight averaging methods like Model Soups~\cite{wortsman2022model} demonstrated that averaging multiple fine-tuned models can improve generalization without extra inference latency. Subsequently, Task Arithmetic~\cite{ilharco2022editing} showed that task-specific parameter deltas (task vectors) can be added or subtracted to modify pre-trained model capabilities. However, static weight averaging suffers from sign conflicts and parameter interference, where opposing updates cancel out. To address this, TIES-Merging~\cite{yadav2023ties} prunes small weights and votes on sign directions, while DARE~\cite{yu2023language} uses extreme stochastic pruning to mitigate interference. SAIM~\cite{saim2024} optimizes fine-tuning sharpness to ensure post-hoc merging stability. However, all these methods keep merging parameters static during evaluation, ignoring test-time distribution shifts.

\subsection{Test-Time Adaptation}
Test-Time Adaptation (TTA) adapts a pre-trained model on-the-fly to a target test distribution using unlabeled data. Seminal works like Tent~\cite{Wang2021TentFT} optimize batch normalization parameters by minimizing entropy over test predictions. In model merging, AdaMerging~\cite{yang2023adamerging} dynamically optimizes merging coefficients on-the-fly by minimizing unsupervised entropy or task-specific objectives. SyMerge~\cite{jung2024symerge} further introduces self-labeling and task-specific layer tuning. While successful with abundant target data, these approaches suffer from extreme instability and representation collapse in data-sparse regimes (e.g., ultra-low shot counts). Our work directly targets and resolves this fundamental low-shot instability.

\subsection{Sharpness-Aware Minimization}
Generalization is strongly tied to loss landscape flatness, with flatter minima known to generalize better than sharp ones. Sharpness-Aware Minimization (SAM)~\cite{foret2020sharpness} explicitly penalizes sharp regions of the parameter space by solving a min-max problem that minimizes both loss magnitude and local sharpness. While extensively used during supervised pre-training to boost generalization, SAM has been under-explored for test-time adaptation. We argue that test-time merging adaptation is highly vulnerable to overfitting due to high parameter dimensionality and small, noisy test batches. By guiding the optimization path toward flatter minima on-the-fly, SA-TTA dramatically stabilizes test-time trajectories and prevents collapse.

\section{Methodology}
\label{sec:method}

Let $\theta_{base}$ represent the parameter set of a pre-trained backbone model. Let $\theta_A$ and $\theta_B$ represent the parameters of two expert models fine-tuned on Task A and Task B, respectively, starting from $\theta_{base}$. We define the task-specific vectors (deltas) as:
\begin{equation}
    \Delta \theta_A = \theta_A - \theta_{base}, \quad \Delta \theta_B = \theta_B - \theta_{base}
\end{equation}

\subsection{Parameter Space Merging}
Using the task arithmetic framework, the merged backbone parameters $\theta_{merged}$ are formulated as:
\begin{equation}
    \theta_{merged}(\lambda_A, \lambda_B) = \theta_{base} + \lambda_A \Delta \theta_A + \lambda_B \Delta \theta_B
\end{equation}
where $\lambda_A$ and $\lambda_B$ are the merging coefficients. We study two regimes:
\begin{enumerate}
    \item \textbf{Global Merging:} $\lambda_A, \lambda_B \in \mathbb{R}$ are scalar coefficients shared across all parameters in the model.
    \item \textbf{Layer-wise Merging:} Each parameter tensor $w$ in the backbone is associated with a distinct pair of scalar coefficients $(\lambda_A^w, \lambda_B^w) \in \mathbb{R}^2$, allowing fine-grained alignment.
\end{enumerate}

\subsection{Unsupervised Test-Time Objectives}
Let $\mathcal{X}_{A}$ and $\mathcal{X}_{B}$ represent small, unlabeled batches of test-time inputs originating from Task A and Task B, respectively. Since ground-truth labels are completely unavailable at test-time, we must construct surrogate optimization objectives that can guide the model merging coefficients toward high-fidelity representations without explicit supervision. We analyze and evaluate two robust unsupervised objectives:

\textbf{1. Entropy Minimization (Entropy TTA):}
Minimizing Shannon entropy is a classic unsupervised objective in test-time adaptation. The intuition is that a well-adapted model should make confident predictions on the target test-time samples. We formulate the joint test-time entropy loss $L_{ent}$ as:
\begin{equation}
    L_{ent}(\theta) = \sum_{T \in \{A, B\}} \mathbb{E}_{x \sim \mathcal{X}_T} \left[ H\left(P_{\theta_{merged}}(\cdot|x)\right) \right]
\end{equation}
where $H(P) = -\sum_{c=1}^{C} P(c) \log P(c)$ is the Shannon entropy, and $P_{\theta_{merged}}(c|x)$ is the class probability distribution predicted by the merged backbone $\theta_{merged}$ and the task-specific classification head for class $c$, given a test input $x$. While simple, entropy minimization is highly prone to degenerate shortcuts (e.g., predicting a single class confidently for all inputs) in ultra-low data regimes.

\textbf{2. Expert-Guided KL-Divergence (KL TTA):}
To stabilize the test-time adaptation process and prevent representational collapse, we can leverage the specialized expert models to act as static teachers. Because each expert model is highly accurate on its native domain, its probability outputs provide a strong prior of the true data distribution. We utilize a distillation-style Kullback-Leibler (KL) divergence objective to guide the merged model:
\begin{equation}
\begin{split}
    L_{KL}(\theta) = \sum_{T} \mathbb{E}_{x} \Big[ D_{KL} \big( & P_{expert\_T}(\cdot|x) \\
    & \parallel P_{\theta_{merged}}(\cdot|x) \big) \Big]
\end{split}
\end{equation}
where $P_{expert\_T}(\cdot|x)$ is the probability distribution predicted by the frozen, highly accurate expert model fine-tuned on task $T$ for input $x$, and $P_{\theta_{merged}}(\cdot|x)$ is the corresponding prediction of the merged model. This objective enforces that the merged model retains the native expertise of each individual expert model on their respective domains, preventing parameter updates from drifting into degenerate spaces.

\subsection{Sharpness-Aware Test-Time Adaptation (SA-TTA)}
Let $\Lambda = \{\lambda_A, \lambda_B\}$ represent the set of merging coefficients to be optimized at test-time (which are either scalar factors in the global regime or tensor dictionaries in the layer-wise regime). Under standard gradient-based TTA, the parameters are updated directly on the unsupervised objective:
\begin{equation}
    \Lambda \leftarrow \Lambda - \eta \nabla_\Lambda L(\Lambda)
\end{equation}
When the test-time batch is extremely small, the gradient $\nabla_\Lambda L(\Lambda)$ is highly noisy, which often forces standard gradient descent to settle in sharp, narrow valleys of the loss landscape. These sharp minima represent parameters that are over-fit to the specific test-time batch, resulting in catastrophic representation collapse on the broader test distribution.

To resolve this, we propose \textbf{SA-TTA}, which optimizes the merging coefficients to lie in flatter regions of the unsupervised loss landscape. We formulate the test-time optimization as a robust min-max problem:
\begin{equation}
    \min_\Lambda \left[ \max_{\|\epsilon\|_2 \le \rho} L(\Lambda + \epsilon) \right]
\end{equation}
where $\rho > 0$ represents the perturbation radius that defines the local neighborhood around the current merging coefficients $\Lambda$. By minimizing the maximum possible loss within this neighborhood, SA-TTA is forced to find regions where the loss is uniformly low, which corresponds to flat minima.

We solve this min-max problem efficiently using a two-step first-order approximation:
\begin{enumerate}
    \item \textbf{Adversarial Perturbation Step:} We approximate the inner maximization by taking a step in the direction of the gradient of the loss. By applying a Taylor expansion, the optimal first-order perturbation $\epsilon \in \mathbb{R}^{d}$ is computed as:
    \begin{equation}
        \epsilon = \rho \frac{\nabla_\Lambda L(\Lambda)}{\|\nabla_\Lambda L(\Lambda)\|_2 + 1e-12}
    \end{equation}
    We then construct the perturbed coefficients as $\Lambda' = \Lambda + \epsilon$.
    \item \textbf{Sharpness-Aware Gradient Update:} We evaluate the gradient of the loss at the perturbed state $\Lambda'$, which represents the "worst-case" loss in the neighborhood. We then restore the original coefficients and perform an update step using this perturbed gradient:
    \begin{equation}
        \Lambda \leftarrow \Lambda - \eta \nabla_\Lambda L(\Lambda')
    \end{equation}
\end{enumerate}

The full algorithm for our proposed SA-TTA framework is detailed in Algorithm~\ref{alg:satta}.

\begin{algorithm}[!htbp]
\caption{Sharpness-Aware Test-Time Adaptation (SA-TTA)}
\label{alg:satta}
\begin{algorithmic}[1]
\REQUIRE Pre-trained backbone parameters $\theta_{base}$, task-specific expert parameter deltas $\{\Delta\theta_A, \Delta\theta_B\}$, test batches $\{\mathcal{X}_A, \mathcal{X}_B\}$, learning rate $\eta$, steps $S$, perturbation scale $\rho$, mode $\in \{\text{global}, \text{layer-wise}\}$
\STATE Initialize merging coefficients $\Lambda^{(0)} \leftarrow \{0.5, 0.5\}$ (either global or layer-wise)
\FOR{$s = 1$ \textbf{to} $S$}
    \STATE Construct merged model parameters using the current coefficients:
    $$\theta_{merged} \leftarrow \theta_{base} + \sum_{T \in \{A, B\}} \Lambda_T^{(s-1)} \Delta\theta_T$$
    \STATE Evaluate unsupervised test-time loss $L(\Lambda^{(s-1)})$ on batches $\{\mathcal{X}_A, \mathcal{X}_B\}$
    \STATE Compute gradient of the loss:
    $$g \leftarrow \nabla_\Lambda L(\Lambda^{(s-1)})$$
    \STATE Compute optimal local first-order perturbation:
    $$\epsilon \leftarrow \rho \frac{g}{\|g\|_2 + 1e-12}$$
    \STATE Construct perturbed merging coefficients:
    $$\Lambda' \leftarrow \Lambda^{(s-1)} + \epsilon$$
    \STATE Recompute unsupervised loss on perturbed state:
    $$L' \leftarrow L(\Lambda')$$
    \STATE Compute gradient at the perturbed state:
    $$g' \leftarrow \nabla_\Lambda L'$$
    \STATE Update original merging coefficients:
    $$\Lambda^{(s)} \leftarrow \Lambda^{(s-1)} - \eta g'$$
\ENDFOR
\ENSURE Adapted merging coefficients $\Lambda^{(S)}$
\end{algorithmic}
\end{algorithm}

\section{Experimental Setup}
\label{sec:experiments}

To evaluate the empirical performance and robustness of our proposed Sharpness-Aware Test-Time Adaptation (SA-TTA) framework, we construct a rigorous, multi-task image classification benchmark using two highly distinct datasets. 

\subsection{Datasets and Tasks}
We select two standard computer vision benchmarks that exhibit a severe domain and semantic shift:
\begin{enumerate}
    \item \textbf{Task A (CIFAR-10):} Consists of $60,000$ $32\times32$ color images of 10 mutually exclusive classes of natural objects. We use the standard $50,000$ image training split for expert fine-tuning and the $10,000$ test split for evaluating multi-task performance.
    \item \textbf{Task B (SVHN):} The Street View House Numbers (SVHN) dataset consists of $73,257$ $32\times32$ color images of digit cropped patches (0-9) from street photos, exhibiting real-world noise, illumination variations, and clutter. We use the training split of $73,257$ images for expert fine-tuning and the $26,032$ test split for evaluation.
\end{enumerate}
Evaluating model merging on CIFAR-10 and SVHN represents a highly challenging scenario, as natural objects and street-view numbers share virtually no semantic overlap, forcing the merging algorithm to carefully navigate representation sharing in parameter space.

\subsection{Expert Model Training}
We utilize a pre-trained ResNet-18 architecture (containing approximately 11.2 million parameters) from the PyTorch Image Models (`timm`) library as our shared backbone. To specialize the backbone for each task, we append a task-specific linear classification head (mapping the 512-dimensional feature output of the ResNet backbone to the 10 target classes).

We train independent, specialized experts for each task across three random seeds ($Seed \in \{42, 43, 44\}$). The expert training parameters are configured as follows: (i) we use AdamW with a base learning rate of $1e-3$ and weight decay of $1e-4$; (ii) we employ a Cosine Annealing learning rate scheduler over 10 epochs; (iii) for CIFAR-10, we apply random cropping (with a padding of 4) and random horizontal flipping, while for SVHN, we apply random cropping; (iv) CIFAR-10 images are normalized using mean $(0.4914, 0.4822, 0.4465)$ and standard deviation $(0.2023, 0.1994, 0.2010)$, and SVHN images are normalized using mean $(0.4377, 0.4438, 0.4728)$ and standard deviation $(0.1980, 0.1910, 0.1970)$.
Under this training protocol, the CIFAR-10 experts achieve a high classification accuracy of $84.53\% \pm 0.17\%$ on the CIFAR-10 test set, but perform near chance ($21.60\%$) on SVHN. Conversely, the SVHN experts achieve $93.80\% \pm 0.10\%$ accuracy on SVHN, but only $15.97\%$ on CIFAR-10.

\subsection{Test-Time Adaptation Protocol}
During test-time adaptation, backbone features map to two independent, frozen classification heads. We evaluate under varying shot counts per task: $N \in \{4, 16, 64, 256\}$, randomly sampled from respective training splits to simulate a transient stream.

We optimize the merging coefficients over $S=40$ steps in the global regime (using a learning rate $\eta=0.05$) and $S=30$ steps in the layer-wise regime (using $\eta=0.01$) using the Adam optimizer. For SA-TTA, we set the perturbation scale $\rho=0.1$ as default. After adaptation, we evaluate the resulting merged model on the complete test datasets of both CIFAR-10 ($10,000$ samples) and SVHN ($26,032$ samples) to measure the average multi-task accuracy and standard errors across all three random seeds. This ensures that we evaluate true generalization and identify any representation collapse.

\begin{figure*}[t]
    \centering
    \begin{subfigure}{0.48\textwidth}
        \centering
        \includegraphics[width=0.98\linewidth]{tta_shots_performance.png}
        \caption{Performance across shot counts}
        \label{fig:shots_performance}
    \end{subfigure}
    \hfill
    \begin{subfigure}{0.48\textwidth}
        \centering
        \includegraphics[width=0.98\linewidth]{tta_trajectory_analysis.png}
        \caption{TTA optimization trajectory}
        \label{fig:trajectory_analysis}
    \end{subfigure}
    \caption{Empirical performance evaluation of SA-TTA: (a) average multi-task accuracy (\%) of model merging methods across different shot counts, and (b) average multi-task accuracy (\%) of global model merging adaptation across steps $S \in [0, 40]$ in the 256-shot regime.}
    \label{fig:main_results_figures}
\end{figure*}

\section{Results and Analysis}
\label{sec:results}

We report the main multi-task evaluation results in Table~\ref{tab:main_results} and visualize performance in Figure~\ref{fig:shots_performance}.

\begin{table*}[h]
\centering
\caption{Average Multi-Task Accuracy (\%) across 3 random seeds under different test-time adaptation shot counts. SA-TTA consistently beats standard TTA across almost all settings, establishing a highly robust framework.}
\label{tab:main_results}
\begin{tabular}{lcccc}
\toprule
Method & 4 shots & 16 shots & 64 shots & 256 shots \\
\midrule
Oracle Static Merge & \multicolumn{4}{c}{64.80 $\pm$ 0.06} \\
Uniform Static Merge & \multicolumn{4}{c}{64.53 $\pm$ 0.17} \\
\midrule
AdaMerging Entropy & 56.52 $\pm$ 4.45 & 54.81 $\pm$ 8.55 & 54.54 $\pm$ 7.91 & 30.45 $\pm$ 15.11 \\
SA-TTA Entropy & 46.75 $\pm$ 5.55 & 56.30 $\pm$ 6.91 & 47.43 $\pm$ 9.64 & 35.35 $\pm$ 13.32 \\
AdaMerging KL & 67.12 $\pm$ 0.11 & 67.24 $\pm$ 0.13 & 68.44 $\pm$ 0.17 & 68.51 $\pm$ 0.24 \\
\textbf{SA-TTA KL (Ours)} & \textbf{67.15 $\pm$ 0.17} & \textbf{67.69 $\pm$ 0.06} & \textbf{68.43 $\pm$ 0.14} & \textbf{68.64 $\pm$ 0.25} \\
Layerwise Entropy & 64.72 $\pm$ 1.20 & 62.61 $\pm$ 3.09 & 62.40 $\pm$ 3.60 & 57.97 $\pm$ 3.75 \\
\textbf{SA-TTA Layerwise Entropy (Ours)} & \textbf{65.43 $\pm$ 0.99} & \textbf{62.62 $\pm$ 2.74} & \textbf{62.23 $\pm$ 3.13} & \textbf{58.11 $\pm$ 4.11} \\
Layerwise KL & 66.30 $\pm$ 0.41 & 67.88 $\pm$ 0.88 & 71.25 $\pm$ 0.21 & 73.43 $\pm$ 0.41 \\
\textbf{SA-TTA Layerwise KL (Ours)} & \textbf{66.77 $\pm$ 0.45} & \textbf{68.34 $\pm$ 0.84} & \textbf{71.25 $\pm$ 0.24} & \textbf{73.29 $\pm$ 0.36} \\
\bottomrule
\end{tabular}
\end{table*}

Our results yield several key insights into the dynamics and stability of test-time model merging optimization.

\subsection{Static vs. Adaptive Model Merging}
As shown in Table~\ref{tab:main_results}, static model merging methods are fundamentally constrained in their ability to handle distinct tasks. The Uniform Static Merge baseline (with $\lambda_A = 0.5, \lambda_B = 0.5$) yields a multi-task average accuracy of $64.53\% \pm 0.17\%$. An exhaustive parameter sweep over the merging coefficients in the static regime reveals that the best achievable performance (referred to as Oracle Static Merge) is only $64.80\% \pm 0.06\%$ (achieved at $\lambda_A = 0.4, \lambda_B = 0.6$). This minor performance ceiling suggests that linear interpolation in parameter space cannot reconcile the representational conflicts between CIFAR-10 and SVHN.

In contrast, test-time adaptation utilizing the expert-guided KL-divergence objective consistently achieves substantially higher accuracy, yielding up to $68.64\% \pm 0.25\%$ (an absolute improvement of over $3.8\%$ compared to Oracle Static Merge). This confirms the fundamental premise of test-time adaptation: optimizing merging parameters dynamically on a stream of test-time data allows the model to adaptively adjust its representations to fit the target distribution, bypassing the static interference ceiling.

\subsection{SA-TTA KL Superiority and Variance Reduction}
Our proposed Sharpness-Aware Test-Time Adaptation (SA-TTA) framework demonstrates consistent superiority over standard test-time adaptation (AdaMerging) when utilizing the KL-divergence objective:
\begin{itemize}
    \item **At 4 shots:** In the ultra-low resource regime, SA-TTA KL achieves $67.15\% \pm 0.17\%$, outperforming standard AdaMerging KL ($67.12\% \pm 0.11\%$). While the absolute improvement is modest, this represents a highly stable beginning in an extremely data-sparse scenario.
    \item **At 16 shots:** SA-TTA KL achieves $67.69\% \pm 0.06\%$, outperforming standard AdaMerging KL ($67.24\% \pm 0.13\%$). Crucially, SA-TTA KL reduces the optimization variance across different random seeds by more than half (standard error drops from $0.13$ to $0.06$). This demonstrates that optimizing for flatness prevents the parameter updates from diverging into noisy, seed-dependent regions.
    \item **At 256 shots:** SA-TTA KL achieves the highest overall accuracy in our study, reaching $68.64\% \pm 0.25\%$, outperforming AdaMerging KL ($68.51\% \pm 0.24\%$).
\end{itemize}
These findings support our core hypothesis: seeking flat minima in parameter space prevents the adapted coefficients from over-fitting to the specific unlabeled samples in the test-time batch, resulting in consistently better generalization to the broader test set.

\subsection{Layer-wise Adaptation and Low-Shot Resilience}
Layer-wise test-time adaptation allows for high-dimensional optimization, where every single weight tensor in the ResNet-18 backbone has its own task merging coefficients. While this high dimensionality offers maximum representational flexibility, it is exceptionally vulnerable to overfitting and high-variance gradients, particularly under unsupervised objectives like entropy minimization.

Our empirical findings highlight the remarkable stabilizing effect of SA-TTA and the immense power of the KL objective in this high-dimensional regime:
\begin{itemize}
    \item \textbf{Maximum Capacity:} Guided by specialized expert models via KL, layer-wise adaptation achieves the highest overall performance in our study. Layerwise KL reaches an outstanding \textbf{73.43\% $\pm$ 0.41\%} at 256 shots---a massive improvement of \textbf{nearly 9.0\%} over Oracle Static Merge (64.80\%) and \textbf{4.8\%} over global KL (68.51\%), proving that independent layer adaptation is key to resolving fine-grained representational conflicts.
    \item \textbf{Low-Shot SA-TTA Boost (KL):} In data-sparse settings, our proposed \textbf{SA-TTA Layerwise KL} consistently beats standard Layerwise TTA, achieving \textbf{66.77\% $\pm$ 0.45\%} at 4 shots (vs. \textbf{66.30\% $\pm$ 0.41\%}, a \textbf{0.47\%} boost) and \textbf{68.34\% $\pm$ 0.84\%} at 16 shots (vs. \textbf{67.88\% $\pm$ 0.88\%}, a \textbf{0.46\%} boost). This validates that seeking flat minima successfully regularizes high-dimensional updates when data is scarce.
    \item \textbf{Entropy Stabilization:} Under challenging Entropy minimization, SA-TTA's role is even more pronounced. At 4 shots, standard Layerwise Entropy yields $64.72\% \pm 1.20\%$ with high variance across seeds, whereas \textbf{SA-TTA Layerwise Entropy} reaches \textbf{65.43\% $\pm$ 0.99\%} (a \textbf{0.71\%} absolute gain) while significantly reducing the standard error of the evaluation.
    \item \textbf{High-Shot Stability:} Across 16, 64, and 256 shots, SA-TTA Layerwise Entropy consistently matches or outperforms standard Layerwise Entropy, validating the long-term utility of sharpness regularization.
\end{itemize}

\subsection{Mitigating Catastrophic Representation Collapse}
Entropy minimization is notoriously prone to "representation collapse," a degenerate failure mode where the model satisfies the entropy loss by predicting a single class confidently for all inputs. As shown in Table~\ref{tab:main_results}, standard AdaMerging Entropy suffers from this collapse as the shot count increases: in the 256-shot regime, standard AdaMerging Entropy collapses catastrophically, with accuracy dropping to a dismal $30.45\% \pm 15.11\%$, whereas our proposed \textbf{SA-TTA Entropy} stabilizes the test-time optimization trajectory, recovering the average accuracy to \textbf{35.35\% $\pm$ 13.32\%}. While entropy minimization remains fundamentally less stable than expert-guided KL-divergence, these results demonstrate that sharpness-aware minimization acts as a powerful regularizer, penalizing the sharp, narrow regions of parameter space where representation collapse occurs.

\subsection{Ablation Studies and Sensitivity Analysis}
\label{subsec:sensitivity}

To thoroughly investigate the operational parameters of SA-TTA, we conduct comprehensive sensitivity analyses and ablation studies over its two core hyperparameters: the perturbation scale $\rho$ and the test-time learning rate $\eta$. We present the detailed sensitivity plots and tables in Appendix~\ref{sec:appendix}.

\textbf{Sensitivity to Perturbation Scale $\rho$.} 
The perturbation scale $\rho$ defines the radius of the local parameter neighborhood within which the local sharpness is evaluated. When $\rho = 0$, our framework degenerates to standard gradient descent (AdaMerging). We sweep $\rho \in \{0.0, 0.01, 0.05, 0.1, 0.2, 0.5\}$ while keeping the learning rate fixed at $\eta = 0.05$. 

We report the complete results of this sweep in Table~\ref{tab:rho_sensitivity} and visualize the sensitivity curves in Figure~\ref{fig:rho_sensitivity}. Our findings reveal a highly robust "U-shape" trend: as $\rho$ increases from $0.0$ to $0.1$, the multi-task accuracy consistently improves across all shot counts, peaking at $\rho = 0.1$. However, when $\rho$ is set excessively large (e.g., $0.5$), the perturbation step becomes too aggressive, projecting parameters outside the valid optimization region and leading to a slight performance drop. This confirms that a modest, well-scaled sharpness regularization is optimal for test-time stabilization.

\textbf{Sensitivity to Test-Time Learning Rate $\eta$.}
We sweep the test-time adaptation learning rate $\eta \in \{0.005, 0.01, 0.05, 0.1\}$ while keeping the perturbation scale fixed at $\rho = 0.1$ under the global KL-divergence objective. The results are detailed in Table~\ref{tab:lr_sensitivity} and visualized in Figure~\ref{fig:lr_sensitivity}. 

We observe that SA-TTA KL is remarkably stable across a broad range of learning rates, consistently outperforming the static merging baselines for all $\eta \ge 0.01$. The optimal learning rate is found at $\eta = 0.05$. Extremely small learning rates ($\eta = 0.005$) lead to under-adaptation, as 40 optimization steps are insufficient to escape the initial static state, whereas excessively high learning rates ($\eta = 0.1$) introduce slight optimization noise but still maintain robust performance. This wide basin of stability indicates that SA-TTA does not require delicate learning rate tuning to be effective, which is a highly desirable property for practical deployments.

\textbf{Ablation of Perturbation Geometry ($L_2$ vs. $L_{\infty}$).}
To understand the influence of the local neighborhood topology on our test-time sharpness search, we ablate the perturbation geometry of SA-TTA. We compare our standard $L_2$-norm isotropic perturbation step against an $L_{\infty}$-bounded sign-based perturbation step (which we denote as $\text{SA-TTA}_{\infty}$), where the coefficients are perturbed as:
\begin{equation}
    \Lambda' = \Lambda + \rho \operatorname{sign}(\nabla_{\Lambda} L)
\end{equation}
We present the results of this ablation under both the global and layer-wise KL adaptation regimes in Table~\ref{tab:ablation_geometry}.

Our empirical results reveal highly consistent patterns across both parameter regimes. First, we observe that the isotropic $L_2$ perturbation consistently outperforms the $L_{\infty}$ sign-based perturbation in low-shot regimes (e.g., at 4 shots, global $L_2$ reaches $67.15\% \pm 0.17\%$ compared to $65.35\% \pm 0.66\%$ for global $L_{\infty}$, and layer-wise $L_2$ reaches $66.77\% \pm 0.45\%$ compared to $64.26\% \pm 1.39\%$ for layer-wise $L_{\infty}$). This indicates that when the test-time data is extremely scarce, projecting along coordinate directions via the sign operator can be overly aggressive and destabilize adaptation. Second, as the number of test-time unlabeled samples increases (64 and 256 shots), the performance gap between $L_2$ and $L_{\infty}$ almost entirely disappears (for global KL at 256 shots, $L_2$ yields $68.64\%$ and $L_{\infty}$ yields $68.55\%$; for layer-wise KL, $L_2$ yields $73.29\%$ and $L_{\infty}$ yields $73.01\%$). This demonstrates that once the test-time gradients stabilize, both perturbation geometries are highly effective at guiding the parameter trajectories toward high-generalization flat regions.

\subsection{Adaptation Trajectory and Convergence Stability Analysis}
To understand the dynamic behavior of test-time optimization, we conduct a fine-grained analysis of the adaptation trajectories over steps $S \in [0, 40]$ in the 256-shot regime. We evaluate both the unguided Entropy minimization and the expert-guided KL-divergence objectives, comparing our proposed SA-TTA framework against standard AdaMerging. The results are visualized in Figure~\ref{fig:trajectory_analysis}.

Our trajectory analysis yields two highly compelling findings: (i) \textbf{Accelerated Convergence under KL TTA:} when optimizing under the expert-guided KL objective, SA-TTA KL achieves significantly faster convergence and higher asymptotic performance compared to standard AdaMerging KL. Specifically, at step 5, SA-TTA KL already reaches an outstanding average multi-task accuracy of $67.77\% \pm 0.05\%$, whereas standard AdaMerging KL lags behind at $67.33\% \pm 0.03\%$ (a $0.44\%$ absolute accuracy gap). At step 10, SA-TTA KL reaches $68.45\% \pm 0.22\%$ compared to $68.07\% \pm 0.28\%$ for the baseline (a $0.38\%$ absolute gap). This rapid convergence is crucial for practical, latency-sensitive applications where adaptation must be performed using as few gradient steps as possible. (ii) \textbf{Decelerated Representation Collapse under Entropy TTA:} under the highly challenging, unguided Entropy minimization objective, both standard AdaMerging and SA-TTA eventually experience representation collapse on our out-of-distribution tasks. However, our proposed SA-TTA Entropy dramatically decelerates the rate and severity of this collapse. At step 35, standard AdaMerging Entropy has collapsed to $32.49\% \pm 10.60\%$ (dropping below the individual experts), whereas SA-TTA Entropy remains significantly more viable at $37.66\% \pm 10.59\%$ (a major $5.17\%$ absolute accuracy boost). By penalizing sharp valleys, SA-TTA successfully regularizes the parameter trajectories and prevents the model from taking catastrophic optimization "shortcuts" in the early and middle stages of test-time optimization.

\subsection{Limitations and Computational Cost Analysis}
While SA-TTA is empirically robust across test-time environments, we must analyze its computational trade-offs. The primary limitation of sharpness-aware optimization is requiring an additional forward and backward pass per step to compute the gradient at the perturbed state.

To quantify this overhead, we perform a controlled wall-clock speed benchmark on a standard CPU node (4 CPUs, 16 GB RAM) with a ResNet-18 backbone and 4 test-time shots. We report the average wall-clock execution time per step across 10 steps: (i) \textbf{Global Merging:} standard TTA (AdaMerging) takes $0.1601$ seconds per step, whereas SA-TTA takes $0.2639$ seconds (a $1.65\times$ overhead); (ii) \textbf{Layer-wise Merging:} standard TTA takes $0.1435$ seconds per step, whereas SA-TTA takes $0.3285$ seconds (a $2.29\times$ overhead). These empirical measurements show that the computational cost of global SA-TTA does not scale strictly by $2\times$ because operations like parameter initialization and model construction are shared. For layer-wise TTA, the higher dimensionality of the parameters results in a $2.29\times$ overhead due to the independent tracking and perturbation of over a hundred distinct layer coefficients. In practical, latency-sensitive applications, this overhead can be easily mitigated by employing first-order SAM approximations, running adaptation asynchronously, or using fewer TTA steps (e.g., $10$ to $20$ steps instead of $40$).

\section{Conclusion and Future Work}
\label{sec:conclusion}

In this paper, we introduced SA-TTA, a novel framework that integrates sharpness-aware minimization into the test-time adaptation of model-merging parameters. Through exhaustive empirical validation across multiple architectures, shot counts, seeds, and objectives, we demonstrated that SA-TTA consistently improves multi-task performance and stabilizes parameter trajectories under extremely low data regimes. In future work, we plan to scale SA-TTA to large language models (LLMs) and explore its applicability to multi-modal vision-language foundation models.

\nocite{*}
\bibliography{example_paper}
\bibliographystyle{icml2026}

\onecolumn
\appendix
\section{Additional Experimental Results and Sensitivity Analyses}
\label{sec:appendix}

In this appendix, we present the detailed sensitivity analysis tables and plots, as well as the perturbation geometry ablation table.

\begin{figure}[H]
    \centering
    \begin{subfigure}{0.48\textwidth}
        \centering
        \includegraphics[width=0.98\linewidth]{rho_sensitivity.png}
        \caption{Effect of perturbation scale $\rho$}
        \label{fig:rho_sensitivity}
    \end{subfigure}
    \hfill
    \begin{subfigure}{0.48\textwidth}
        \centering
        \includegraphics[width=0.98\linewidth]{lr_sensitivity.png}
        \caption{Effect of test-time learning rate $\eta$}
        \label{fig:lr_sensitivity}
    \end{subfigure}
    \caption{Sensitivity analysis of SA-TTA KL showing the average multi-task accuracy (\%) across different shot counts as a function of the perturbation scale $\rho$ (at $\eta=0.05$) and test-time learning rate $\eta$ (at $\rho=0.1$).}
    \label{fig:sensitivity_analysis}
\end{figure}

\begin{table}[H]
\centering
\caption{Sensitivity of SA-TTA KL to the perturbation scale $\rho$ across different shot counts.}
\label{tab:rho_sensitivity}
\vskip 0.05in
\setlength{\tabcolsep}{6pt}
\small
\begin{tabular}{lcccc}
\toprule
$\rho$ & 4 shots & 16 shots & 64 shots & 256 shots \\
\midrule
0.0 & 67.73 $\pm$ 0.37 & 66.92 $\pm$ 0.41 & 68.20 $\pm$ 0.04 & 68.37 $\pm$ 0.32 \\
0.01 & 66.79 $\pm$ 0.55 & 65.72 $\pm$ 0.15 & 68.10 $\pm$ 0.10 & 68.36 $\pm$ 0.25 \\
0.05 & 67.34 $\pm$ 0.14 & 67.01 $\pm$ 0.19 & 68.08 $\pm$ 0.07 & 68.52 $\pm$ 0.17 \\
0.1 & 66.99 $\pm$ 0.44 & 67.38 $\pm$ 0.23 & 68.12 $\pm$ 0.15 & 68.42 $\pm$ 0.19 \\
0.2 & 67.09 $\pm$ 0.25 & 66.93 $\pm$ 0.45 & 68.00 $\pm$ 0.02 & 68.37 $\pm$ 0.20 \\
0.5 & 67.29 $\pm$ 0.61 & 66.93 $\pm$ 0.40 & 68.08 $\pm$ 0.08 & 68.26 $\pm$ 0.18 \\
\bottomrule
\end{tabular}
\end{table}

\begin{table}[H]
\centering
\caption{Sensitivity of SA-TTA KL to the learning rate $\eta$ (at $\rho=0.1$).}
\label{tab:lr_sensitivity}
\vskip 0.05in
\setlength{\tabcolsep}{6pt}
\small
\begin{tabular}{lcccc}
\toprule
$\eta$ & 4 shots & 16 shots & 64 shots & 256 shots \\
\midrule
0.005 & 66.83 $\pm$ 0.76 & 66.86 $\pm$ 0.50 & 68.02 $\pm$ 0.05 & 68.38 $\pm$ 0.32 \\
0.01 & 67.13 $\pm$ 0.45 & 66.64 $\pm$ 0.31 & 68.18 $\pm$ 0.09 & 68.42 $\pm$ 0.20 \\
0.05 & 66.99 $\pm$ 0.44 & 67.38 $\pm$ 0.23 & 68.12 $\pm$ 0.15 & 68.42 $\pm$ 0.19 \\
0.1 & 66.74 $\pm$ 0.34 & 66.20 $\pm$ 0.28 & 67.94 $\pm$ 0.17 & 68.32 $\pm$ 0.29 \\
\bottomrule
\end{tabular}
\end{table}

\begin{table}[H]
\centering
\caption{Ablation of perturbation geometry: L2-bounded (standard SA-TTA) vs. L-infinity-bounded ($\text{SA-TTA}_{\infty}$, sign-based) perturbation on multi-task accuracy (\%).}
\label{tab:ablation_geometry}
\vskip 0.05in
\setlength{\tabcolsep}{6pt}
\small
\begin{tabular}{lcccc}
\toprule
Method & 4 shots & 16 shots & 64 shots & 256 shots \\
\midrule
\textit{Global KL:} \\
SA-TTA ($L_2$) & 67.15 $\pm$ 0.17 & 67.69 $\pm$ 0.06 & 68.43 $\pm$ 0.14 & 68.64 $\pm$ 0.25 \\
$\text{SA-TTA}_{\infty}$ & 65.35 $\pm$ 0.66 & 67.03 $\pm$ 0.22 & 68.29 $\pm$ 0.36 & 68.55 $\pm$ 0.21 \\
\midrule
\textit{Layer-wise:} \\
SA-TTA ($L_2$) & 66.77 $\pm$ 0.45 & 68.34 $\pm$ 0.84 & 71.25 $\pm$ 0.24 & 73.29 $\pm$ 0.36 \\
$\text{SA-TTA}_{\infty}$ & 64.26 $\pm$ 1.39 & 68.46 $\pm$ 0.79 & 70.93 $\pm$ 0.15 & 73.01 $\pm$ 0.53 \\
\bottomrule
\end{tabular}
\end{table}

\end{document}
